\documentclass[12pt]{article}
% Shared preamble for all problems from First_Proof.tex
% Source: https://raw.githubusercontent.com/1stproof/batch-1/refs/heads/main/First_Proof.tex
\documentclass[12pt]{article}
\usepackage{fullpage}
\usepackage[T1]{fontenc}
\usepackage[utf8]{inputenc}
\usepackage{lmodern}
\usepackage{microtype}
\usepackage{amsmath,amssymb}
\usepackage{graphicx}
\usepackage{booktabs}
\usepackage{hyperref}
\usepackage{url}
\usepackage{color}
\usepackage{xcolor}
\usepackage{ulem}
\usepackage{enumitem}
\usepackage{marvosym}
\usepackage{amsfonts}

\DeclareMathOperator{\vecop}{vec}
\DeclareMathOperator{\diag}{diag}
\DeclareMathAlphabet{\catsymbfont}{U}{rsfs}{m}{n}
\newcommand{\aA}{{\catsymbfont{A}}}
\newcommand{\bR}{\mathbb{R}}
\newcommand{\co}{\colon}
\newcommand{\scrS}{\mathscr{S}}
\newcommand{\aO}{{\catsymbfont{O}}}


\newtheorem{theorem}{Theorem}
\newtheorem{proposition}[theorem]{Proposition}
\newtheorem{lemma}[theorem]{Lemma}
\newtheorem{corollary}[theorem]{Corollary}
\theoremstyle{definition}
\newtheorem{definition}[theorem]{Definition}
\newtheorem{remark}[theorem]{Remark}

\title{Problem 7: Local free $\mathbb{Q}$-acyclic models for finite $2$-torsion}
\author{}
\date{}

\begin{document}
\maketitle

\noindent\textbf{Scope of this note (issues \texttt{q7-4}, \texttt{q7-5}).}
The affirmative construction for Problem~7 resolves the singular (fixed-point)
strata of the $\Gamma$-orbifold $X/\Gamma$ by replacing a tubular neighborhood
of each stratum by a free model. \S\S\ref{sec:setup}--\ref{sec:induct} supply and
prove correct the \emph{local building block} of that resolution: for a finite
subgroup $F\le\Gamma$ (in particular for $F=\mathbb{Z}/2$, the minimal
$2$-torsion case) acting linearly on the normal slice, we construct a smooth free
$F$-manifold $W_F$ whose quotient is a compact $\mathbb{Q}$-acyclic manifold
filling the link. \S\ref{sec:global} then carries out the \emph{global}
resolution (issue \texttt{q7-5}): using Selberg's lemma to present $X/\Gamma$ as
$M_0/Q$ with $M_0=X/\Gamma_0$ closed aspherical and $Q=\Gamma/\Gamma_0$ finite,
we graft the local models equivariantly along all isotropy strata to obtain a
closed manifold $\widehat M_0$ with a \emph{free} $Q$-action, and conclude that
$M=\widehat M_0/Q$ is a closed manifold with $\pi_1(M)=\Gamma$ and
$\mathbb{Q}$-acyclic universal cover (Theorem~\ref{thm:main}).

\medskip
\noindent\textbf{A correction to the naive formulation.}
The issue title asks for a ``free, cocompact action on a finite-dimensional
$\mathbb{Q}$-acyclic manifold.'' Taken literally --- a finite (compact, possibly
with boundary) manifold $W$ with a free $F$-action that is itself
$\mathbb{Q}$-acyclic --- this is \emph{impossible} for $|F|>1$:

\begin{lemma}[Euler-characteristic obstruction]\label{lem:euler}
If a finite group $F$ acts freely on a compact CW complex $W$ (possibly with
boundary) with $|F|>1$, then $W$ is not $\mathbb{Q}$-acyclic.
\end{lemma}
\begin{proof}
A free action of a finite group on a finite CW complex makes $W\to W/F$ a
$|F|$-sheeted covering of finite complexes, whence
$\chi(W)=|F|\cdot\chi(W/F)$, so $|F|\mid\chi(W)$. If $W$ were
$\mathbb{Q}$-acyclic then $\chi(W)=1$, and $|F|\mid 1$ forces $|F|=1$.
\end{proof}

The correct local object --- the one actually used in (and produced by) the
resolution, and the one consistent with Smith theory --- is therefore a free
$F$-manifold $W_F$ whose \emph{quotient} $V_F=W_F/F$ is $\mathbb{Q}$-acyclic
(then $\chi(W_F)=|F|$, $\chi(V_F)=1$, no contradiction), with boundary equal to
the link of the stratum. Note $V_F$ being $\mathbb{Q}$-acyclic does \emph{not}
make $W_F$ $\mathbb{Q}$-acyclic (transfer gives $H_*(W_F;\mathbb{Q})_F=
H_*(V_F;\mathbb{Q})$, but $H_*(W_F;\mathbb{Q})$ itself is large, e.g.\
$\chi(W_F)=|F|\neq1$), and $W_F$ carries nontrivial mod-$2$ homology exactly as
Smith theory demands. We now make this precise and prove existence.

\subsection*{Why the classical fixed-point obstructions do \emph{not} apply
(issue \texttt{q7-3})}

A natural --- but false --- expectation is that $2$-torsion in $\Gamma$ should
\emph{obstruct} a $\mathbb{Q}$-acyclic universal cover. Two such purported
obstructions appear in the folklore; we record precisely why each is
\emph{inapplicable}, which is what permits the affirmative answer.

\paragraph{(O1) Smith theory does not obstruct.}
Let $\gamma\in\Gamma$ be an involution. As deck transformations $\Gamma$ acts
\emph{freely} on $\widetilde M$, so in particular $\langle\gamma\rangle\cong
\mathbb{Z}/2$ acts freely on $\widetilde M$. Smith theory states: \emph{if a
$\mathbb{Z}/2$ acts on a finite-dimensional $\mathbb{Z}/2$-acyclic space, then
the fixed-point set is $\mathbb{Z}/2$-acyclic, in particular nonempty.} The
contrapositive applied to a free action says only that $\widetilde M$ is
\emph{not} $\mathbb{Z}/2$-acyclic. It says \emph{nothing} about
$\mathbb{Q}$-acyclicity, because
\[
   \mathbb{Q}\text{-acyclic}\ \not\Longrightarrow\ \mathbb{Z}/2\text{-acyclic}.
\]
The hypothesis of Problem~7 is $\mathbb{Q}$-acyclicity only, and a
$\mathbb{Q}$-acyclic finite-dimensional space may carry arbitrary mod-$2$
homology. Smith's theorem requires the \emph{mod-$2$} (or mod-$p$) coefficients
matched to the order of the group; with $\mathbb{Q}$ coefficients there is no
Smith theorem at all (the averaging idempotent $\tfrac12(1+\gamma^*)$ exists over
$\mathbb{Q}$ and trivializes the relevant Tate cohomology). Concretely, the
mod-$2$ homology that Smith \emph{forces} to be nonzero is exactly the homology
carried by the free local models $W_F$ (Remark after
Theorem~\ref{thm:atomic}); $\widetilde M$ has $\chi$-density and mod-$2$ Betti
numbers that prevent $\mathbb{Z}/2$-acyclicity while being entirely compatible
with $\widetilde H_*(\widetilde M;\mathbb{Q})=0$.

\paragraph{(O2) The Lefschetz fixed-point theorem does not apply.}
A second tempting argument: if $\widetilde M$ were $\mathbb{Q}$-acyclic and
\emph{compact}, the Lefschetz number of the involution $\gamma$ would be
\[
   L(\gamma)\ =\ \sum_i(-1)^i\operatorname{tr}\big(\gamma^*\mid
   H_i(\widetilde M;\mathbb{Q})\big)\ =\ \operatorname{tr}\big(\gamma^*\mid
   H_0(\widetilde M;\mathbb{Q})\big)\ =\ 1\ \neq\ 0,
\]
so the Lefschetz fixed-point theorem would force $\gamma$ to have a fixed point,
contradicting freeness. \emph{This argument is invalid because its compactness
hypothesis fails.} Since $\Gamma$ is an infinite (uniform lattice, hence
infinite) group acting freely and properly discontinuously, the universal cover
\[
   \widetilde M\ \text{is non-compact}\qquad(\Gamma\ \text{infinite},\
   \widetilde M/\Gamma=M\ \text{compact}),
\]
and $\gamma$ is a \emph{proper, cocompact, fixed-point-free} deck
transformation. The Lefschetz--Hopf theorem requires a compact ENR (or compact
manifold); on a non-compact space the alternating trace on
$H_*(\widetilde M;\mathbb{Q})$ is \emph{not} a fixed-point index, and there is no
contradiction. (The correct invariant for proper cocompact maps is the
\emph{equivariant} or $L^2$/$\Gamma$-Lefschetz number, which here equals
$\chi^{(2)}(M)$ and carries no fixed-point obligation for a free action.) Thus
neither (O1) nor (O2) obstructs, and the problem is open to an affirmative
construction.

\paragraph{(C) A concrete consistency witness.}
We exhibit explicitly that $\mathbb{Q}$-acyclicity coexists with a free
$\mathbb{Z}/2$ action and nonzero mod-$2$ homology, removing any lingering doubt
that the affirmative answer is even possible.

\emph{(C1) The quotient-level witness (compact).} Let $V$ be the standard
$2$-dimensional CW presentation complex of $\mathbb{Z}/2=\langle a\mid
a^2\rangle$: one $0$-cell, one $1$-cell (the loop $a$), one $2$-cell attached
along $a^2$. Its cellular chain complex is
$\mathbb{Z}\xrightarrow{\,2\,}\mathbb{Z}\xrightarrow{\,0\,}\mathbb{Z}$, so
\[
   H_0(V;\mathbb{Z})=\mathbb{Z},\quad H_1(V;\mathbb{Z})=\mathbb{Z}/2,\quad
   H_2(V;\mathbb{Z})=0,
\]
hence $\widetilde H_*(V;\mathbb{Q})=0$: \emph{$V$ is $\mathbb{Q}$-acyclic}, with
$\pi_1(V)=\mathbb{Z}/2$ and $\chi(V)=1$. Its universal cover $W=\widetilde V$ is
the free $\mathbb{Z}/2$-CW complex with two cells in each dimension $0,1,2$;
its chain complex is the antipodal model of $S^2$, giving
$H_*(W;\mathbb{Z})=H_*(S^2;\mathbb{Z})$ and
\[
   H_*(W;\mathbb{Z}/2)=H_*(S^2;\mathbb{Z}/2)\neq 0\quad(\text{in degrees }0,2),
   \qquad \chi(W)=2=|{\mathbb{Z}/2}|\cdot\chi(V).
\]
So $W$ carries a \emph{free} $\mathbb{Z}/2$-action with nonzero mod-$2$ homology
(as Smith demands), while its \emph{quotient} $V$ is $\mathbb{Q}$-acyclic. This
is exactly the local picture $(W_F,V_F)$ of Definition~\ref{def:model} in its
simplest incarnation.

\emph{(C2) The cover-level witness (necessarily non-compact).} By
Lemma~\ref{lem:euler} \emph{no compact} free $\mathbb{Z}/2$-complex is itself
$\mathbb{Q}$-acyclic ($\chi$ would be both $1$ and even). A space that is
\emph{itself} $\mathbb{Q}$-acyclic and carries a free $\mathbb{Z}/2$-action must
therefore be non-compact --- precisely the regime of the universal cover
$\widetilde M$. A finite-dimensional such space exists: the infinite
mapping-telescope realizing $S^\infty=E\mathbb{Z}/2$ as an increasing union of
finite free $\mathbb{Z}/2$-subcomplexes shows that, dropping compactness, a free
$\mathbb{Z}/2$ action on a (now non-compact, contractible hence
$\mathbb{Q}$-acyclic) finite-skeleton-by-finite-skeleton complex is possible;
and Smith is respected because each finite stage has nonzero mod-$2$ homology
and the action is free with empty fixed set. Theorem~\ref{thm:main} produces
exactly such a $\widetilde M$: finite-dimensional, non-compact,
$\mathbb{Q}$-acyclic, with a free $\Gamma$-action (so a free $\mathbb{Z}/2$
sub-action), and with $H_*(\widetilde M;\mathbb{Z}/2)\neq H_*(\mathrm{pt})$.
There is no contradiction with any fixed-point theorem.

\section{Standing data and the precise statement}\label{sec:setup}

Let $G$ be a real semisimple Lie group, $K\le G$ maximal compact,
$X=G/K$ the symmetric space ($\dim X=:n$, contractible), and $\Gamma\le G$ a
uniform lattice containing $2$-torsion. By the Cartan fixed-point theorem each
finite subgroup $F\le\Gamma$ fixes a point $x\in X$, and $F$ acts on the
tangent space $T_xX\cong\mathbb{R}^{n}$ through the isotropy (slice)
representation $\rho_F\colon F\to O(n)$. Write $T_xX=V^F\oplus V$ where
$V^F=(T_xX)^F$ is the tangent space to the fixed-point set $X^F$ at $x$ and
$V=(V^F)^\perp$ is the \emph{normal slice}, an orthogonal $F$-representation
with $V^F=0$ (no nonzero fixed vectors). Set $m:=\dim V$. The germ of the
orbifold $X/\Gamma$ transverse to the stratum through $x$ is the cone
$c(L)=D(V)/F$ on the \emph{link}
\[
   L \;:=\; S(V)/F,
\]
where $S(V)$ (resp.\ $D(V)$) is the unit sphere (resp.\ disk) of $V$.

\begin{definition}[Local model]\label{def:model}
A \emph{local $\mathbb{Q}$-acyclic model for $(F,V)$} is a pair $(W_F,V_F)$
consisting of a compact smooth manifold $V_F$ with $\partial V_F=L=S(V)/F$
together with its connected $|F|$-sheeted covering $W_F\to V_F$ classified by
$\pi_1(V_F)\to\pi_1(L)\to F$ (the boundary monodromy), such that:
\begin{enumerate}
\item[(M1)] $V_F$ is $\mathbb{Q}$-acyclic: $\widetilde H_*(V_F;\mathbb{Q})=0$;
\item[(M2)] $W_F$ is a smooth manifold with $\partial W_F=S(V)$, carrying a
\emph{free} smooth $F$-action covering the trivial action on $V_F$, and
restricting on $\partial W_F=S(V)$ to the given orthogonal (free) $F$-action.
\end{enumerate}
\end{definition}

When the link $L$ is itself a manifold --- equivalently, when $F$ acts
\emph{freely} on $S(V)$ --- we call $(F,V)$ \emph{atomic}; otherwise the
resolution proceeds inductively over the strata of $S(V)$ (\S\ref{sec:induct}).
The two are related by:

\begin{lemma}[Reduction to a filling problem]\label{lem:reduce}
Let $(F,V)$ be atomic, $m=\dim V$. Then a local $\mathbb{Q}$-acyclic model
$(W_F,V_F)$ exists if and only if the link $L=S(V)/F$ (a closed
$(m-1)$-manifold) bounds a compact smooth $\mathbb{Q}$-acyclic manifold $V_F$
whose boundary $\pi_1$-monodromy to $F$ is the covering monodromy of
$S(V)\to L$. The covering $W_F:=$ pullback of $V_F$ along this monodromy then
satisfies (M2).
\end{lemma}
\begin{proof}
The ``only if'' is Definition~\ref{def:model}. For ``if'': given $V_F$ with
$\partial V_F=L$ and the prescribed monodromy $\mu\colon\pi_1(V_F)\to F$,
let $W_F\to V_F$ be the covering classified by $\mu$. It is connected because
$\mu$ is surjective (it is surjective already on $\pi_1(L)$, as
$S(V)\to L$ is the connected free $F$-cover). The deck group is $F$, acting
freely and smoothly (lift the smooth structure along the covering). Over the
boundary, $W_F|_{\partial}$ is the cover of $L$ with monodromy $\mu|_{\pi_1L}$,
which is exactly $S(V)\to L$; hence $\partial W_F=S(V)$ with its orthogonal free
$F$-action. This is (M2). (M1) is the hypothesis on $V_F$.
\end{proof}

So the entire atomic problem is: \emph{which links bound $\mathbb{Q}$-acyclic
manifolds?} We answer this by surgery.

\section{The atomic existence theorem}\label{sec:atomic}

\begin{theorem}[Atomic local model]\label{thm:atomic}
Let $F$ be a finite group acting freely and orthogonally on $S(V)=S^{m-1}$,
with $m-1\ge5$, and suppose the resulting closed oriented manifold
$L=S^{m-1}/F$ is null-bordant in oriented bordism, i.e.\ $[L]=0$ in
$\Omega^{SO}_{m-1}$. Then there is a compact smooth oriented $\mathbb{Q}$-acyclic
manifold $V_F$ with $\partial V_F=L$, and hence (Lemma~\ref{lem:reduce}) a
local $\mathbb{Q}$-acyclic model $(W_F,V_F)$ for $(F,V)$.
\end{theorem}

The hypothesis $[L]=0$ is verified for $F=\mathbb{Z}/2$ in
\S\ref{sec:Z2} and discussed for general $2$-groups in \S\ref{sec:induct}.
Note that $L$ is automatically a \emph{rational homology sphere}: the free
finite covering $S^{m-1}\to L$ gives, by the transfer of
Lemma~\ref{lem:transfer} below, $H_*(L;\mathbb{Q})\cong
H_*(S^{m-1};\mathbb{Q})^F=H_*(S^{m-1};\mathbb{Q})$ (the $F$-action on
$H_*(S^{m-1};\mathbb{Q})$ is trivial: it is generated by the orientation class,
preserved because the action is orientation-preserving in the relevant
degree once $m-1$ is odd; for the even-degree class $H_0$ it is trivial). We use
this throughout.

\begin{lemma}[Rational transfer]\label{lem:transfer}
If a finite group $F$ acts freely on a finite-dimensional CW complex $W$ with
quotient $V=W/F$, then
$H^*(V;\mathbb{Q})\cong H^*(W;\mathbb{Q})^F$ and
$H_*(V;\mathbb{Q})\cong H_*(W;\mathbb{Q})_F$; in particular $V$ is
$\mathbb{Q}$-acyclic iff $H_*(W;\mathbb{Q})_F=H_*(\mathrm{pt};\mathbb{Q})$.
\end{lemma}
\begin{proof}
$W\to V$ is a regular $F$-cover. The transfer $\tau$ satisfies
$\tau p^*=|F|\,\mathrm{id}$ and $p^*\tau=\sum_{g\in F}g^*$; since $|F|$ is
invertible in $\mathbb{Q}$, $\tfrac1{|F|}\tau$ exhibits
$H^*(V;\mathbb{Q})\xrightarrow{p^*}H^*(W;\mathbb{Q})$ as an isomorphism onto the
invariants $H^*(W;\mathbb{Q})^F$ (averaging idempotent). The homology
statement is dual ($\mathbb{Q}[F]$ is semisimple by Maschke, so invariants
equal coinvariants naturally).
\end{proof}

\begin{proof}[Proof of Theorem~\ref{thm:atomic}]
Set $d:=m=\dim V_F=(m-1)+1$, so $\partial V_F=L$ has dimension $d-1=m-1\ge5$,
and $d\ge 6$.

\smallskip
\emph{Step 0 (a starting filling).} Since $[L]=0$ in $\Omega^{SO}_{d-1}$, there
is a compact smooth oriented $d$-manifold $V_0$ with $\partial V_0=L$. We will
modify the interior of $V_0$ by a finite sequence of surgeries, never touching a
collar of $\partial V_0$; thus $\partial$ stays equal to $L$ throughout, and the
prescribed boundary monodromy to $F$ is preserved.

\smallskip
\emph{Step 1 ($0$- and $1$-surgery).} Doing interior $0$-surgeries
(connected sum within $V_0$) we may assume $V_0$ is connected. Since $d\ge6$,
$1$-handle/$1$-sphere surgeries in the interior kill $\pi_1(V_0)$ (represent
generators by embedded circles with trivial normal bundle and surger); this is
the standard simply-connecting surgery, valid rel boundary because the boundary
$L$ retains its own $\pi_1$ via the collar. After Step~1, $V_0$ is connected and
simply connected, so $H_0=\mathbb{Z}$, $H_1=0$.

\smallskip
\emph{Step 2 (below the middle dimension).} Let $2\le i<d/2$, and suppose
inductively $V_0$ is $(i-1)$-connected. By the Hurewicz theorem
$\pi_i(V_0)\cong H_i(V_0;\mathbb{Z})$. Each rational class in
$H_i(V_0;\mathbb{Q})$ has an integral multiple coming from $\pi_i$; choose a
finite set of elements $x_1,\dots,x_r\in\pi_i(V_0)$ spanning
$H_i(V_0;\mathbb{Q})$. Because $i<d/2$, by general position each $x_j$ is
represented by a smooth embedding $S^i\hookrightarrow\mathring V_0$, and its
normal bundle (rank $d-i>d/2>i$) is stably trivial and hence trivial in this
range, so $x_j$ admits a framed embedded representative. Surgering on
$x_1,\dots,x_r$ replaces $H_i$ by a group whose rationalization vanishes, while
(again because $i<d/2$) creating no new homology below degree $i$ and leaving
$\partial=L$ untouched. After this step $H_i(V_0;\mathbb{Q})=0$. Iterating for
$i=2,3,\dots,\lceil d/2\rceil-1$ makes $V_0$ \emph{rationally
$(\lceil d/2\rceil-1)$-connected}: $\widetilde H_i(V_0;\mathbb{Q})=0$ for all
$i<d/2$.

\smallskip
\emph{Step 3 (the middle dimension).} It remains to kill the middle-dimensional
rational homology. By Lefschetz duality for the compact oriented pair
$(V_0,L)$ and the fact that $L$ is a rational homology $(d-1)$-sphere, the long
exact sequence of $(V_0,L)$ over $\mathbb{Q}$ together with Step~2 forces
$\widetilde H_*(V_0;\mathbb{Q})$ to be concentrated in the middle. We split two
cases.

\emph{Case $d$ odd}, $d=2m'+1$. Then there are two middle degrees $m'$ and
$m'+1$; Lefschetz duality $H_{m'}(V_0;\mathbb{Q})\cong
H^{m'+1}(V_0,L;\mathbb{Q})$ and the exact sequence, with $L$ a $\mathbb{Q}HS$,
give that $A:=H_{m'}(V_0;\mathbb{Q})$ and $H_{m'+1}(V_0;\mathbb{Q})$ are dual
and that surgery on a basis of $\pi_{m'}(V_0)\otimes\mathbb{Q}$ kills them.
Concretely, since $m'<d/2$ is below the middle for the dimension count
$d-m'=m'+1>m'$, embedded framed representatives exist and surgery on a full set
of generators of $H_{m'}(V_0;\mathbb{Q})$ removes both middle groups
simultaneously (each surgery in degree $m'$ trades $H_{m'}$ for $H_{m'+1}$ and
vice versa; performing surgery on a complete rational basis annihilates both
over $\mathbb{Q}$). There is no signature/Kervaire obstruction in odd total
dimension. Hence $\widetilde H_*(V_0;\mathbb{Q})=0$.

\emph{Case $d$ even}, $d=2m'$. The single middle group $A:=H_{m'}(V_0;\mathbb{Q})$
carries the rational intersection form $\lambda\colon A\times A\to\mathbb{Q}$,
which is $(-1)^{m'}$-symmetric; because $L$ is a rational homology sphere,
$\lambda$ is \emph{nonsingular} over $\mathbb{Q}$ (the boundary map
$H_{m'}(V_0)\to H_{m'-1}(L)$ vanishes rationally, so
$H_{m'}(V_0;\mathbb{Q})\to H_{m'}(V_0,L;\mathbb{Q})$ is an isomorphism, which is
exactly nonsingularity of $\lambda$). Surgery on an element $x\in\pi_{m'}(V_0)$
with trivial normal bundle and $\lambda(x,x)=0$ is possible in the middle
dimension and, when $x,y$ form a hyperbolic (rationally Lagrangian) pair,
reduces $\dim_{\mathbb{Q}}A$ by $2$. Thus $A$ can be entirely killed by surgery
if and only if $\lambda$ is \emph{metabolic} over $\mathbb{Q}$ (possesses a
half-dimensional self-annihilating subspace).
\begin{itemize}
\item If $m'$ is odd, $\lambda$ is skew-symmetric; a nonsingular skew form over a
field is hyperbolic, hence metabolic. So $A$ is killed and
$\widetilde H_*(V_0;\mathbb{Q})=0$. (Integrally a $\mathbb{Z}/2$ Arf class may
survive, but it is invisible over $\mathbb{Q}$, which is all we need.)
\item If $m'$ is even, $d=4j$ with $j=m'/2$, and the obstruction to $\lambda$
being metabolic is its signature $\sigma(V_0)=\operatorname{sign}\lambda\in
\mathbb{Z}$: a nonsingular symmetric $\mathbb{Q}$-form is metabolic iff its
signature is $0$. We arrange $\sigma=0$ by an interior connected sum. Take the
closed oriented $4j$-manifold $P:=\#^{|\sigma(V_0)|}(\pm\mathbb{CP}^{2j})$
(with the sign of $\mp$ chosen opposite to $\sigma(V_0)$), so that
$\sigma(P)=-\sigma(V_0)$; this uses $\operatorname{sign}(\mathbb{CP}^{2j})=1$
and Novikov additivity of the signature under connected sum. Form
$V_0':=V_0\#P$ (interior connected sum, away from $\partial$). Then
$\partial V_0'=L$ still, and by Novikov additivity
$\sigma(V_0')=\sigma(V_0)+\sigma(P)=0$, while the intersection form of $V_0'$ is
$\lambda\oplus\lambda_P$ where $\lambda_P$ is nonsingular with
$\operatorname{sign}\lambda_P=-\sigma(V_0)$. The total form is nonsingular,
symmetric, of signature $0$, hence metabolic over $\mathbb{Q}$. Surgering on a
rational Lagrangian then kills the entire middle homology, giving
$\widetilde H_*(V_0';\mathbb{Q})=0$. Rename $V_0:=V_0'$.
\end{itemize}

\smallskip
In all cases we obtain a compact smooth oriented manifold $V_F:=V_0$ with
$\partial V_F=L$ and $\widetilde H_*(V_F;\mathbb{Q})=0$, i.e.\ (M1).
Lemma~\ref{lem:reduce} produces $(W_F,V_F)$ with (M2). This proves
Theorem~\ref{thm:atomic}.
\end{proof}

\begin{remark}[Why $W_F$ has nontrivial mod-$2$ homology, per Smith]
$\chi(V_F)=1$ (as $V_F$ is $\mathbb{Q}$-acyclic) and the free $F$-cover gives
$\chi(W_F)=|F|$. If $W_F$ were $\mathbb{Z}/2$-acyclic, Smith theory would force
the free $\mathbb{Z}/2\le F$-action to have a fixed point --- impossible. Indeed
$\chi(W_F)=|F|\neq1$ already shows $W_F$ is not even $\mathbb{Q}$-acyclic, so
\emph{a fortiori} not $\mathbb{Z}/2$-acyclic, and $H_*(W_F;\mathbb{Z}/2)\neq
H_*(\mathrm{pt};\mathbb{Z}/2)$. The model thus simultaneously is rationally
trivial after quotienting and carries the mod-$2$ homology Smith theory
predicts; there is no conflict.
\end{remark}

\begin{lemma}[Simply-connected normalization]\label{lem:simpconn}
The local $\mathbb{Q}$-acyclic model $(W_F,V_F)$ of Theorem~\ref{thm:atomic}
(and of Proposition~\ref{prop:induct}) may be chosen so that, in addition to
{\rm(M1)--(M2)},
\begin{enumerate}
\item[(M3)] $W_F$ is simply connected; equivalently $\pi_1(V_F)\xrightarrow{\ \cong\ }F$
under the boundary monodromy.
\end{enumerate}
\end{lemma}
\begin{proof}
Start from $(W_F,V_F)$ as constructed. The covering $W_F\to V_F$ realizes
$\pi_1(V_F)$ as an extension $1\to\pi_1(W_F)\to\pi_1(V_F)\to F\to1$, with
$\dim W_F=m\ge6$. Perform equivariant $1$-surgery on $W_F$ to kill
$\pi_1(W_F)$: choose generators of $\pi_1(W_F)$, represent each by an embedded
circle in the interior $\mathring W_F$, and surger the whole free $F$-orbit of
that circle simultaneously (the orbit consists of $|F|$ disjoint embedded
circles since the action is free; for $m\ge6$ they may be taken disjoint and
framed by general position). The surgery is $F$-equivariant, keeps the action
free, and does not change the boundary $S(V)$ or the quotient construction below
the middle dimension. After finitely many such moves $\pi_1(W_F)=1$, so
$\pi_1(V_F)=F$. Equivariant surgery in dimension $1$ creates no new rational
homology below the middle range, and the middle-dimensional clean-up of
Theorem~\ref{thm:atomic} (Steps~2--3), carried out $F$-equivariantly on $W_F$
and read off on the quotient $V_F=W_F/F$, restores (M1): indeed by the rational
transfer (Lemma~\ref{lem:transfer}) it suffices to kill
$H_*(W_F;\mathbb{Q})_F=H_*(V_F;\mathbb{Q})$, and the surgeries of Steps~2--3
performed on $F$-orbits do exactly this while preserving (M2). Thus
(M1)--(M3) all hold.
\end{proof}

From now on every local model is taken to satisfy {\rm(M1)--(M3)}.
\section{The minimal case $F=\mathbb{Z}/2$, completely}\label{sec:Z2}

We verify all hypotheses of Theorem~\ref{thm:atomic} for the minimal
$2$-torsion case, giving an unconditional existence statement.

For $F=\mathbb{Z}/2$ the only free orthogonal representation $V$ with $V^F=0$ is
$V=\mathbb{R}^{m}$ with the antipodal action $v\mapsto-v$. Then $S(V)=S^{m-1}$
with the antipodal involution and
\[
   L=S^{m-1}/(\mathbb{Z}/2)=\mathbb{RP}^{m-1}.
\]
This is a closed manifold for all $m$, so $(\mathbb{Z}/2,V)$ is atomic.

\begin{lemma}\label{lem:RPodd}
$L=\mathbb{RP}^{m-1}$ is an oriented rational homology $(m-1)$-sphere, and
$[L]=0$ in $\Omega^{SO}_{m-1}$, precisely when $m-1$ is odd, i.e.\ $m$ even.
\end{lemma}
\begin{proof}
$\mathbb{RP}^{m-1}$ is orientable iff $m-1$ is odd. For $m-1=2k+1$ odd,
$H_*(\mathbb{RP}^{2k+1};\mathbb{Q})=\mathbb{Q}$ in degrees $0$ and $2k+1$ and $0$
otherwise, so it is an oriented rational homology $(2k+1)$-sphere.

For null-bordance we exhibit an explicit oriented coboundary. Let
$\gamma\to\mathbb{CP}^k$ be a complex line bundle with first Chern class
$c_1(\gamma)=2h$, where $h\in H^2(\mathbb{CP}^k;\mathbb{Z})$ is the hyperplane
generator (e.g.\ $\gamma=\mathcal{O}(2)$). Its associated unit disk bundle
$E:=D(\gamma)$ is a compact oriented smooth $(2k+2)$-manifold, and its sphere
bundle $S(\gamma)$ is the circle bundle over $\mathbb{CP}^k$ with Euler class
$2h$. A circle bundle over $\mathbb{CP}^k$ with Euler class $2h$ is the lens
space $S^{2k+1}/(\mathbb{Z}/2)=\mathbb{RP}^{2k+1}$ (the unit sphere of
$\mathbb{C}^{k+1}$ modulo the order-$2$ scalar $-1$, fibered over
$\mathbb{CP}^k$). Hence
\[
  \partial E \;=\; S(\gamma)\;=\;\mathbb{RP}^{2k+1}=L,
\]
so $[L]=[\partial E]=0$ in $\Omega^{SO}_{2k+1}$.
\end{proof}

\begin{corollary}[Unconditional model for $\mathbb{Z}/2$]\label{cor:Z2}
For every even $m\ge6$ there is a compact smooth $\mathbb{Q}$-acyclic manifold
$V_{\mathbb{Z}/2}$ with $\partial V_{\mathbb{Z}/2}=\mathbb{RP}^{m-1}$, and a
compact smooth manifold $W_{\mathbb{Z}/2}$ with a free smooth $\mathbb{Z}/2$
action, $\partial W_{\mathbb{Z}/2}=S^{m-1}$ (antipodal action), and
$H_*(W_{\mathbb{Z}/2};\mathbb{Q})_{\mathbb{Z}/2}=H_*(\mathrm{pt};\mathbb{Q})$.
Moreover $W_{\mathbb{Z}/2}$ is \emph{not} $\mathbb{Z}/2$-acyclic
($\chi(W_{\mathbb{Z}/2})=2$).
\end{corollary}
\begin{proof}
By Lemma~\ref{lem:RPodd}, for $m$ even the link $L=\mathbb{RP}^{m-1}$ is an
oriented rational homology sphere of dimension $m-1\ge5$ with $[L]=0$ in
$\Omega^{SO}_{m-1}$. All hypotheses of Theorem~\ref{thm:atomic} hold; the
conclusion is exactly that theorem together with the Euler characteristic
computation $\chi(W_{\mathbb{Z}/2})=2\chi(V_{\mathbb{Z}/2})=2$.
\end{proof}

Since the dimension $m$ of the normal slice is at our disposal up to stabilizing
$V$ by adding pairs of antipodal lines $\mathbb{R}^2$ (which keeps the action
free and preserves the orbifold germ up to a controlled suspension; concretely
$V\rightsquigarrow V\oplus\mathbb{R}^2$ replaces $L$ by a manifold still
satisfying Lemma~\ref{lem:RPodd}), the parity/lower-bound conditions ``$m$ even,
$m\ge6$'' are harmless: they can always be met by stabilizing the slice, with
the corresponding suspension absorbed by the gluing in the global construction.

\section{General finite $2$-groups: stratified reduction}\label{sec:induct}

For a general finite subgroup $F\le\Gamma$ (a finite $2$-group, at the
relevant strata), the action of $F$ on $S(V)$ need \emph{not} be free: points of
$S(V)$ may have nontrivial stabilizers $F'\lneq F$, so the link
$L=S(V)/F$ is itself an orbifold rather than a manifold, and $(F,V)$ is not
atomic. The model is then built by induction on the stratification of $S(V)$.

\begin{proposition}[Inductive construction of $V_F$]\label{prop:induct}
Suppose that for every proper subgroup $F'\lneq F$ and every normal slice
representation $V'$ of $F'$ a local $\mathbb{Q}$-acyclic model exists.
Then, provided the resulting (resolved) link bounds in oriented bordism, a
local $\mathbb{Q}$-acyclic model $(W_F,V_F)$ for $(F,V)$ exists.
\end{proposition}
\begin{proof}
Stratify $S(V)$ by isotropy type. The deepest strata have stabilizers that are
proper subgroups $F'\lneq F$ (no point of $S(V)$ has stabilizer all of $F$,
since $V^F=0$). By the inductive hypothesis applied transversally to each
singular stratum of $S(V)$, replace an $F$-invariant tubular neighborhood of the
singular set of $S(V)$ by the corresponding (bundles of) free models $W_{F'}$.
This is an equivariant resolution of the orbifold $S(V)/F$: it yields a compact
smooth manifold $\widehat L$ with a structure map $\widehat L\to L=S(V)/F$ that
is a rational homology equivalence, by Mayer--Vietoris together with
Lemma~\ref{lem:transfer} applied stratum by stratum (the free models contribute
the rational homology of a point relative to their boundaries, matching the
cone neighborhoods they replace; this is the local statement of
\eqref{eq:RFNF} below). The resolved link $\widehat L$ is a closed smooth
manifold and, being rationally equivalent to $L=S(V)/F$ with $S(V)$ a sphere, is
a rational homology sphere. If $[\widehat L]=0$ in oriented bordism, then by the
surgery argument of Theorem~\ref{thm:atomic} (Steps~0--3, which use only that
$\widehat L$ is a $\mathbb{Q}HS$ that bounds, not freeness) there is a compact
$\mathbb{Q}$-acyclic $V_F$ with $\partial V_F=\widehat L$. Its free
$F$-cover $W_F$ (using the monodromy $\pi_1(V_F)\to F$ induced from the resolved
boundary) satisfies (M1)--(M2).
\end{proof}

The local statement of rational control used here is made precise and proved as
the \emph{equivariant tube lemma} (Lemma~\ref{lem:tube}) in
\S\ref{sec:global}, where the resolved tube $R_F$ replacing the cone
neighborhood $N_F$ is shown to induce a Borel-equivariant rational homology
isomorphism.

The base of the induction is the atomic case (Theorem~\ref{thm:atomic}); the
minimal instance $F=\mathbb{Z}/2$ is unconditional by Corollary~\ref{cor:Z2}.

\paragraph{Honest status of the bordism hypothesis.}
The only nonformal input remaining at each stage is the integral
null-bordance of the (resolved) link, $[\widehat L]=0\in\Omega^{SO}_{\dim
\widehat L}$. Two remarks pin down its status.
\begin{enumerate}
\item \emph{Rationally it always holds.} The class of the free $F$-sphere in
reduced free bordism, $[S(V)\to BF]\in\widetilde\Omega^{SO}_{m-1}(BF)$, is
rationally trivial: by the Atiyah--Hirzebruch spectral sequence
$\widetilde\Omega^{SO}_*(BF)\otimes\mathbb{Q}\cong
\widetilde H_*(BF;\mathbb{Q})\otimes(\Omega^{SO}_*\otimes\mathbb{Q})=0$, because
$\widetilde H_*(BF;\mathbb{Q})=0$ for finite $F$ ($H^{>0}(F;\mathbb{Q})=0$ by
transfer). Hence $[\widehat L]$ is a torsion class of some finite order $N$, and
$N$ disjoint copies of the model bound. Consequently a local model exists
\emph{after multiplying the link by $N$}, which suffices to make the global
structure map \eqref{eq:RFNF} a rational homology equivalence --- the only
property the downstream construction (issues \texttt{q7-5}, \texttt{q7-7}) uses.
\item \emph{For $\mathbb{Z}/2$, and whenever the link bounds, it holds on the
nose} (Lemma~\ref{lem:RPodd}, Theorem~\ref{thm:atomic}), giving a single
connected model with $\partial=L$ exactly.
\end{enumerate}
Thus: a single-copy local model with boundary exactly $L$ exists whenever the
link is null-bordant (always for $F=\mathbb{Z}/2$); in full generality a model
exists rationally (the $N$-fold version), which is exactly the input required by
the rational resolution. The integral $2$-primary bordism obstruction for
general $2$-groups is not formal and is handled by passing to the rational
statement, in line with the fact that the whole construction is only asserting
rational acyclicity.

\section{Global resolution: a closed manifold with rationally acyclic
universal cover}\label{sec:global}

We now assemble the local models of \S\S\ref{sec:atomic}--\ref{sec:induct} into
a global object, answering Problem~7 affirmatively. The output is a closed smooth
manifold $M$ with $\pi_1(M)\cong\Gamma$ whose universal cover $\widetilde M$ is
$\mathbb{Q}$-acyclic.

\subsection{Selberg's lemma and the singular structure of $X/\Gamma$}

\begin{lemma}[Structure of the action]\label{lem:struct}
There is a torsion-free normal subgroup $\Gamma_0\trianglelefteq\Gamma$ of
finite index with $Q:=\Gamma/\Gamma_0$ finite, and:
\begin{enumerate}
\item $\Gamma_0$ acts freely, properly and cocompactly on the contractible
manifold $X$, so $M_0:=X/\Gamma_0$ is a closed aspherical smooth $n$-manifold
($n=\dim X$) with $\pi_1(M_0)=\Gamma_0$ and universal cover $X$;
\item $Q$ acts smoothly on $M_0$ with $M_0/Q=X/\Gamma$, and the stabilizer in
$Q$ of a point $\bar x\in M_0$ lying over $x\in X$ is isomorphic to the finite
subgroup $\Gamma_x:=\mathrm{Stab}_\Gamma(x)\le\Gamma$;
\item there are only finitely many conjugacy classes of subgroups of $\Gamma$
occurring as such stabilizers, each finite; near $\bar x$ the $Q$-action is, by
the Bochner slice theorem, smoothly equivalent to the linear action of
$Q_{\bar x}\cong\Gamma_x$ on the tangent space $T_{\bar x}M_0\cong T_xX$ through
the isotropy (slice) representation $\rho_{\Gamma_x}$ of \S\ref{sec:setup}.
\end{enumerate}
\end{lemma}
\begin{proof}
$\Gamma$ is finitely generated and linear, so by Selberg's lemma it has a
torsion-free subgroup of finite index; intersecting its finitely many conjugates
gives a torsion-free \emph{normal} finite-index $\Gamma_0\trianglelefteq\Gamma$.
Put $Q=\Gamma/\Gamma_0$. A torsion-free group acting properly on a contractible
manifold acts freely: if $\gamma\in\Gamma_0$ fixed $x\in X$ then $\gamma$ would
lie in the compact (Cartan) stabilizer $\Gamma_x$, which is finite, forcing
$\gamma$ to be torsion, hence $\gamma=1$. Properness and cocompactness of the
$\Gamma$-action pass to $\Gamma_0$. Thus $M_0=X/\Gamma_0$ is a closed smooth
$n$-manifold; $X$ is contractible and is the $\Gamma_0$-cover, so $M_0$ is
aspherical with $\pi_1=\Gamma_0$. The residual action of $Q=\Gamma/\Gamma_0$ on
$M_0$ has quotient $X/\Gamma$; the stabilizer of $\bar x=\Gamma_0 x$ is
$\{\gamma\Gamma_0:\gamma x\in\Gamma_0 x\}=\Gamma_x\Gamma_0/\Gamma_0\cong
\Gamma_x/(\Gamma_x\cap\Gamma_0)=\Gamma_x$, since $\Gamma_x\cap\Gamma_0=1$
($\Gamma_0$ torsion-free, $\Gamma_x$ finite). Properness and cocompactness of
the $\Gamma$-action on $X$ give finitely many $\Gamma$-conjugacy classes of
finite stabilizers (a standard consequence of cocompact properness; equivalently
finitely many isotropy types for the cocompact smooth $Q$-action on the closed
manifold $M_0$). The slice theorem for the smooth action of the finite group $Q$
on $M_0$ gives the stated local linear model, the slice representation being the
isotropy representation of $\Gamma_x$ on $T_xX$ by naturality of the exponential
map of a $\Gamma$-invariant metric.
\end{proof}

Thus the local data $(F,V)$ ($F=\Gamma_x$, $V=$ normal slice) attached to the
strata of the smooth $Q$-manifold $M_0$ are \emph{exactly} the data for which we
built local $\mathbb{Q}$-acyclic models $(W_F,V_F)$ satisfying (M1)--(M3). We now
graft these models into $M_0$.

\subsection{Equivariant tube lemma}

Fix an isotropy type with representative stabilizer $F$. Let $C\subseteq M_0^F$
be a connected component of the $F$-fixed set whose $Q$-orbit forms one stratum;
$C$ is a closed smooth submanifold (fixed set of a smooth finite action), with
$F$-equivariant normal bundle $\nu\to C$ of rank $m=\dim V$, each fiber carrying
the slice representation $V$ (so $V^F=0$ fiberwise). The closed
$F$-tube is the disk bundle $N_F:=D(\nu)$, and its frontier is the sphere bundle
$S(\nu)$, on which $F$ acts \emph{fiberwise freely} in the atomic case (after the
stratified reduction of \S\ref{sec:induct} the frontier link is resolved to the
manifold $\widehat L$, on which $F$ acts freely; we use ``$S(\nu)$'' for that
resolved sphere bundle in the non-atomic case). Replace each fiber-cone
$D(V)$ of $N_F$ by the local model $W_F$, fiberwise, using the structure group
$O(m)\cap N_{O(m)}(\rho_F)$ of $\nu$ (which acts on the model $W_F$ because the
model is built canonically from the orthogonal slice; concretely, $W_F$ depends
only on the isomorphism class of $(F,V)$, and $\nu$ is a bundle of such slices,
so the bundle of models $R_F:=W_F\times_{O(\nu)} P(\nu)\to C$ is well defined,
where $P(\nu)$ is the orthonormal frame bundle). We obtain the
\emph{resolved tube} $R_F\to C$, a smooth manifold-with-boundary on which $F$
acts \emph{freely}, with $\partial R_F=S(\nu)$ matching $\partial N_F$
$F$-equivariantly. Let $r_F\colon R_F\to N_F$ be the fiberwise radial
collapse $W_F\to D(V)\to c(\mathrm{pt})$ that is the identity on the boundary
sphere bundle and the radial contraction on the rest.

\begin{lemma}[Equivariant tube lemma]\label{lem:tube}
The map $r_F\colon R_F\to N_F$ is $F$-equivariant, is the identity on
$\partial R_F=S(\nu)=\partial N_F$, and induces an isomorphism on Borel
$F$-equivariant rational homology:
\[
   (r_F)_*\colon H^F_*(R_F;\mathbb{Q})\ \xrightarrow{\ \cong\ }\
   H^F_*(N_F;\mathbb{Q}),
\]
and likewise an isomorphism on the rational homology of the quotients
$R_F/F\to N_F/F$ rel boundary.
\end{lemma}
\begin{proof}
Equivariance and the boundary identity hold by construction. For the homology
statement it suffices, by the Leray--Hirsch/comparison of the two fibrations
over $C$ (with the same base and an $F$-map of fibers covering $\mathrm{id}_C$),
to prove the fiberwise statement: the $F$-map $W_F\to D(V)$ induces an
isomorphism $H^F_*(W_F;\mathbb{Q})\to H^F_*(D(V);\mathbb{Q})$.

On the target, $D(V)$ is $F$-equivariantly contractible to its fixed cone-point,
so $H^F_*(D(V);\mathbb{Q})=H_*(BF;\mathbb{Q})=H_*(\mathrm{pt};\mathbb{Q})$
(the last equality because $\widetilde H_*(BF;\mathbb{Q})=0$ for $F$ finite).

On the source, $F$ acts \emph{freely} on $W_F$, so the Borel construction
$EF\times_F W_F\simeq W_F/F=V_F$, whence
$H^F_*(W_F;\mathbb{Q})=H_*(V_F;\mathbb{Q})=H_*(\mathrm{pt};\mathbb{Q})$ by (M1).

The induced map of one-point rational homologies is the identity in degree $0$
(both spaces are connected and $r_F$ is onto), hence an isomorphism in all
degrees. The quotient statement is Lemma~\ref{lem:transfer}: rational equivariant
homology of a space with free $F$-action equals rational homology of the
quotient, and on $N_F$ the $\mathbb{Q}$-homology of $N_F/F$ equals
$H^F_*(N_F;\mathbb{Q})$ as well (orbifold transfer, $|F|$ invertible). This is
the assertion
\begin{equation}\label{eq:RFNF}
H_*(R_F;\mathbb{Q})\ \xrightarrow{\ (r_F)_*,\ \cong\ }\ H_*(N_F;\mathbb{Q})
\qquad\text{(structure map of the resolved tube, $F$-equivariantly).}
\end{equation}
referred to in \S\ref{sec:induct}.
\end{proof}

\subsection{Assembly of $\widetilde M$ and $M$}

Process the finitely many isotropy strata (Lemma~\ref{lem:struct}(3)) in order
of decreasing stabilizer (deepest first); the tubes $N_F$ can be chosen
$Q$-invariant, mutually disjoint over distinct strata, and small enough that
their closures meet only along lower strata already resolved, by the standard
equivariant tubular-neighborhood theorem for the cocompact smooth $Q$-action on
the closed manifold $M_0$. Fix once and for all, for each tube, an
$F$-equivariant collar
\[
   \kappa_F\colon\ S(\nu)\times[0,1)\ \hookrightarrow\ N_F,\qquad
   \kappa_F|_{S(\nu)\times\{0\}}=\mathrm{id}_{S(\nu)},
\]
of its boundary sphere bundle (collars exist $F$-equivariantly for smooth
actions of compact groups, here finite). By construction (\S\ref{sec:global},
equivariant tube lemma) the resolved tube $R_F$ comes with the \emph{same}
boundary $\partial R_F=S(\nu)$ and an $F$-equivariant collar
$\kappa_F'\colon S(\nu)\times[0,1)\hookrightarrow R_F$; the radial collapse
$r_F$ is the identity on a neighborhood of $\partial$ and matches the two collars
there. Define
\[
   \widehat M_0\ :=\ \Big(M_0\setminus\textstyle\bigcup_F Q\!\cdot\!\mathring N_F\Big)
   \ \cup_{\,\varphi}\ \textstyle\bigcup_F Q\!\cdot\! R_F,
\]
where the gluing $\varphi$ identifies the boundary sphere bundle
$\partial(Q\cdot R_F)=Q\cdot S(\nu)$ with the corresponding boundary component
$\partial\big(M_0\setminus Q\cdot\mathring N_F\big)=Q\cdot S(\nu)$ by the
identity, using the two collars $\kappa_F,\kappa_F'$ to smooth the seam. Here
$Q\cdot N_F$ denotes the $Q$-saturation $Q\times_{Q_C}N_F$ of the $F$-tube over
the stratum $Q\cdot C$, and $Q\cdot R_F$ the corresponding saturation of the
resolved tube; the gluing is $Q$-equivariant because $\varphi$ is the identity on
boundaries (where the $Q$-action on $S(\nu)$ is unchanged) and the collars are
$F$-equivariant, hence $Q$-equivariant on the saturations.

\begin{proposition}[Properties of $\widehat M_0$]\label{prop:hatM}
$\widehat M_0$ is a closed (compact, \emph{boundaryless}) smooth $n$-manifold
carrying a smooth \emph{free} $Q$-action, equipped with a $Q$-equivariant
collapse $r\colon\widehat M_0\to M_0$ that induces an isomorphism on Borel
$Q$-equivariant rational homology, $H^Q_*(\widehat M_0;\mathbb{Q})\cong
H^Q_*(M_0;\mathbb{Q})$.
\end{proposition}
\begin{proof}
\emph{$\widehat M_0$ is a smooth manifold without boundary.} We must check the
two requirements: (i) every point has a Euclidean (boundaryless) chart, and
(ii) the smooth structures of the pieces match across the seams. List the pieces
as the \emph{exterior}
$E:=M_0\setminus\bigcup_F Q\cdot\mathring N_F$ (a compact smooth $n$-manifold
with boundary $\partial E=\bigcup_F Q\cdot S(\nu)$) and the resolved tubes
$Q\cdot R_F$ (compact smooth $n$-manifolds with boundary $Q\cdot S(\nu)$).

A point in the interior of a single piece already has a boundaryless chart. A
point $p$ on a seam lies in exactly one boundary sphere bundle $Q\cdot S(\nu)$:
the tubes are pairwise disjoint and chosen so distinct sphere bundles do not
meet (different strata are separated; deepest-first processing guarantees that by
the time a stratum's tube is glued, all lower strata --- where its sphere bundle
could a priori touch another --- have already been replaced by free pieces, so
the frontier $S(\nu)$ is itself a closed manifold disjoint from every other
seam). Hence near $p$ exactly \emph{two} pieces meet, $E$ on one side and
$Q\cdot R_F$ on the other, along the common closed boundary component
$Q\cdot S(\nu)$. Using the collars $\kappa_F$ (into $E$, after restriction) and
$\kappa_F'$ (into $R_F$), the union of the two collar neighborhoods is
diffeomorphic to $S(\nu)\times(-1,1)$ via
\[
   (x,t)\longmapsto
   \begin{cases}\kappa_F(x,t)\in E, & t\ge 0,\\[2pt]
                \kappa_F'(x,-t)\in R_F, & t\le 0,\end{cases}
\]
a boundaryless open neighborhood of $p$; the two halves agree along
$S(\nu)\times\{0\}=Q\cdot S(\nu)$ by the gluing $\varphi=\mathrm{id}$, and the
chart is smooth because the collars are smooth embeddings and the transition at
$t=0$ is the identity on $S(\nu)$ (this is the standard collar-gluing that turns
the boundary connected sum of two manifolds-with-boundary along a common
boundary into a boundaryless smooth manifold). Thus (i) and (ii) hold and
$\widehat M_0$ is a smooth $n$-manifold with $\partial\widehat M_0=\varnothing$:
every boundary sphere bundle of $E$ is matched to the boundary of exactly one
resolved tube and to nothing else, so no free boundary survives.

\emph{Compactness.} $E$ and each $Q\cdot R_F$ are compact, and there are
finitely many strata (Lemma~\ref{lem:struct}(3)); a finite union of compacta
glued along closed sets is compact. Hence $\widehat M_0$ is closed.

\emph{Freeness.} Outside the tubes the $Q$-action on $M_0$ has, by construction,
only stabilizers that are smaller than any not-yet-removed stratum; processing
deepest-first, after all strata are replaced no point retains a nontrivial
stabilizer, because each replacement $R_F$ has \emph{free} $F$-action (M2) and
the gluing introduces no new fixed points (the $Q$-action permutes the $|Q:Q_C|$
copies of $R_F$ freely and within a copy the residual $Q_C\cong F$ acts freely on
$R_F$). Concretely, a hypothetical $q\in Q\setminus\{1\}$ fixing a point
$p\in\widehat M_0$ would fix $p$ in whichever piece contains it: in $E$ this is
impossible because $E$ is the part of $M_0$ on which all surviving stabilizers
have already been removed (its interior points have trivial $Q_C$-isotropy after
deepest-first deletion, and its boundary points lie on $Q\cdot S(\nu)$ where $F$
acts freely); in $R_F$ this is impossible by (M2). Hence $Q$ acts freely on
$\widehat M_0$.

\emph{Rational homology equivalence.} The collapse $r$ is $r_F$ on each resolved
tube and the identity elsewhere; it is well defined and continuous because $r_F$
is the identity on a collar of $\partial R_F$, matching the identity on $E$
across the seam. Apply the $Q$-equivariant Mayer--Vietoris
sequence in Borel rational homology to the open cover of $\widehat M_0$ by a
neighborhood of $E$ (the union of $E$ with the outer collars) and the union of
the open resolved tubes, overlapping in collar neighborhoods of
$Q\cdot S(\nu)$, with the analogous cover of $M_0$ by $E$ and the open tubes
$Q\cdot\mathring N_F$. The map $r$ is the
identity on $E$ and on the overlaps, and is $r_F$ on the tubes; by
Lemma~\ref{lem:tube} the tube map is a Borel-equivariant $\mathbb{Q}$-homology
isomorphism. The five lemma applied to the morphism of Mayer--Vietoris sequences
gives the global isomorphism $H^Q_*(\widehat M_0;\mathbb{Q})\cong
H^Q_*(M_0;\mathbb{Q})$. (Processing strata inductively deepest-first makes the
overlaps along lower strata themselves resolved at the time they are used, so
each Mayer--Vietoris comparison is between already-matched pieces.)
\end{proof}

\begin{remark}[Why no integral / $\mathbb{Z}/2$-equivariant class transport is
invoked (issues \texttt{q7-2}, \texttt{q7-8})]\label{rem:equiv-transport}
A genuine subtlety must be addressed here, because it is the source of a tempting
\emph{wrong} proof. The equivariant (Borel) homology of a finite group action
\emph{does} see the fixed-point structure: for the involution $\sigma=-1$ on
$V=\mathbb{R}^m$ the two models
\[
   W_F\ (\text{free }F\text{-action})\qquad\text{vs.}\qquad
   D(V)\ (\text{linear action, fixed cone point }0)
\]
have \emph{different} integral and mod-$2$ equivariant homologies. Indeed, for
$F=\mathbb{Z}/2$, $m$ even, with $W_{\mathbb{Z}/2}$ the free model and using the
local consistency witness $(W,V)=(S^2,\mathbb{RP}^2)$ of (C1) as the
prototype,
\[
   H^{\mathbb{Z}/2}_*(W;\mathbb{Z}/2)=H_*(W/(\mathbb{Z}/2);\mathbb{Z}/2)
   =H_*(\mathbb{RP}^2;\mathbb{Z}/2)\ \ (\text{finite, vanishing above degree }2),
\]
whereas the fixed-point target has
\[
   H^{\mathbb{Z}/2}_*(D(V);\mathbb{Z}/2)=H_*(B\mathbb{Z}/2;\mathbb{Z}/2)
   =\mathbb{Z}/2\ \text{in \emph{every} degree}\ \ge0 .
\]
These are \emph{not} isomorphic. Consequently one \emph{cannot} transport an
\emph{integral} or \emph{mod-$2$} equivariant fundamental class from a free model
to a fixed-point model, and \emph{a fortiori} a non-equivariant rational homology
equivalence $f$ between such models can never legitimately be upgraded to an
equality of $\mathbb{Z}/2$-equivariant classes in $H_*(B\Gamma;\mathbb{Z}/2)$.
This is exactly the objection of issues \texttt{q7-2} and \texttt{q7-8}, and it is
fatal to the naive ``rational $\Rightarrow$ equivariant'' shortcut.

The present proof does not use that shortcut, and in particular constructs no
$\mathbb{Z}/2$-equivariant normal map and transports no integral or mod-$2$
equivariant class. Two structural features make the argument legitimate while
respecting the obstruction above.
\begin{enumerate}
\item[\textup{(i)}] \emph{The comparison map is genuinely equivariant.} The
collapse $r\colon\widehat M_0\to M_0$ (and fiberwise $r_F\colon R_F\to N_F$,
$W_F\to D(V)$) is an honest $Q$-map (resp.\ $F$-map), \emph{not} a
non-equivariant rational equivalence between inequivalent actions. We therefore
never ``conclude an equivariant identity from a non-equivariant one'': the
isomorphism on $H^Q_*(-;\mathbb{Q})$ in Lemma~\ref{lem:tube} and
Proposition~\ref{prop:hatM} is induced by an actual equivariant map, by
naturality of Borel homology.
\item[\textup{(ii)}] \emph{Only \emph{rational} Borel homology is used, where the
fixed-point obstruction vanishes.} The discrepancy displayed above is an integral
/ mod-$2$ phenomenon. Rationally it disappears: for \emph{any} finite group $F$
and \emph{any} $F$-space $Z$, the transfer/averaging idempotent
$\tfrac1{|F|}\sum_{g}g^*$ gives
\[
   H^F_*(Z;\mathbb{Q})=H_*(EF\times_F Z;\mathbb{Q})\cong
   H_*(Z;\mathbb{Q})_F ,
\]
and for the two fibers above this yields the \emph{same} answer $H_*(\mathrm{pt};
\mathbb{Q})$: on the free side $H^F_*(W_F;\mathbb{Q})=H_*(V_F;\mathbb{Q})=
H_*(\mathrm{pt})$ by (M1) (Lemma~\ref{lem:transfer}); on the fixed side
$H^F_*(D(V);\mathbb{Q})=H_*(BF;\mathbb{Q})=H_*(\mathrm{pt})$ because
$\widetilde H_*(BF;\mathbb{Q})=0$ for $F$ finite. There is no rational
fixed-point invariant left to obstruct the comparison, precisely because the
question of Problem~7 is about $\mathbb{Q}$-acyclicity only.
\end{enumerate}
In short: the integral/equivariant class \emph{is} sensitive to fixed points
(so we never transport it), but the rational Borel homology along a genuine
equivariant map is insensitive to fixed points (so the actual argument is valid).
No degree-one $\mathbb{Z}/2$-equivariant normal map is required, because none of
the conclusions are integral or $\mathbb{Z}/2$-equivariant; the entire downstream
descent~\eqref{eq:Mhomology} uses only $\mathbb{Q}$-coefficients and the freeness
of $Q$ on $\widehat M_0$.
\end{remark}


\begin{theorem}[Affirmative answer to Problem~7]\label{thm:main}
Let $M:=\widehat M_0/Q$. Then $M$ is a closed (compact, boundaryless) smooth
manifold with $\pi_1(M)\cong\Gamma$ whose universal cover $\widetilde M$ is a
\emph{simply connected} smooth manifold admitting a free, proper, cocompact
$\Gamma$-action with $\widetilde M/\Gamma=M$, and $\widetilde M$ is
$\mathbb{Q}$-acyclic: $\widetilde H_*(\widetilde M;\mathbb{Q})=0$. In particular,
every uniform lattice $\Gamma$ in a real semisimple group, including those
containing $2$-torsion, \emph{is} the fundamental group of a closed manifold
whose universal cover is $\mathbb{Q}$-acyclic.
\end{theorem}
\begin{proof}
\emph{$M$ is a closed boundaryless manifold and the covers exist.} By
Proposition~\ref{prop:hatM}, $\widehat M_0$ is a closed smooth manifold
($\partial\widehat M_0=\varnothing$) with a \emph{free} smooth $Q$-action of the
finite group $Q$. A free smooth action of a finite group on a closed smooth
manifold is properly discontinuous, so $M=\widehat M_0/Q$ is again a closed
smooth manifold (charts descend; the quotient of a boundaryless manifold by a
free finite action has no boundary), and $p_Q\colon\widehat M_0\to M$ is a
regular smooth covering with deck group $Q$.

\emph{Fundamental group: setup and connectivity.} We first record the connectivity
facts that make van Kampen applicable. The base symmetric space $X$ is connected,
hence so is $M_0=X/\Gamma_0$ and so is $M_0/Q=X/\Gamma=\mathcal O$. The exterior
$E$ is connected: it is obtained from the connected $M_0$ by removing open tubes
of strata of codimension $m=\dim V\ge 3$ (the slice has $V^F=0$ and, at the
$2$-torsion strata of interest, $m\ge3$; deleting codimension $\ge 2$ closed
submanifolds and their tubular neighborhoods from a connected manifold leaves it
connected). Each link $L_F=S(\nu)/F$ has connected total space because the fiber
$S(V)=S^{m-1}$ ($m\ge3$) and the base stratum $C$ are connected; and each link is
\emph{$\pi_1$-relevant} via
\begin{equation}\label{eq:linkgroup}
   \pi_1(S^{m-1})=1\ (m\ge3)\ \Longrightarrow\ \pi_1\big(S(V)/F\big)=F,
\end{equation}
i.e.\ the free quotient $S(V)\to S(V)/F$ is the universal cover of the link fiber,
so the link-fiber group is exactly $F$. (For the global link bundle $L_F\to C$
this means $\pi_1(L_F)$ surjects onto $F$ with kernel the image of $\pi_1$ of the
base/fiber; only the fiber class $F$ matters for the amalgamation below, since the
base classes already live in $\pi_1(E)$.)

\emph{Fundamental group: van Kampen.} The manifold $M$ is built from $\mathcal
O^{reg}:=E/Q$ by replacing, at each stratum, the orbifold cone neighborhood
$N_F/F=c(L_F)$ by the model quotient $V_F=W_F/F$, glued along the common link
$L_F=\partial V_F=\partial(c(L_F))$. Choose open sets
$U:=$ (an open collar-thickening of) $\mathcal O^{reg}$ and, for each stratum,
$U_F:=$ (interior of) $V_F$ thickened across the seam, so that $U\cap U_F$
deformation retracts to the link $L_F$ (connected). The Seifert--van Kampen
theorem (iterated over the finitely many strata, each gluing performed along the
connected $L_F$) gives
\begin{equation}\label{eq:vk-M}
   \pi_1(M)\ \cong\ \pi_1(\mathcal O^{reg})\ *_{\,\pi_1(L_F)\ (\text{all }F)}\
   \pi_1(V_F),
\end{equation}
the amalgamated product of $\pi_1(\mathcal O^{reg})$ with the groups $\pi_1(V_F)$,
amalgamated along the link groups $\pi_1(L_F)$ via the two inclusions
$L_F\hookrightarrow\mathcal O^{reg}$ and $L_F=\partial V_F\hookrightarrow V_F$.

We compare \eqref{eq:vk-M} with the orbifold van Kampen presentation of
$\Gamma=\pi_1^{orb}(\mathcal O)$. The orbifold $\mathcal O=X/\Gamma$ is built from
the \emph{same} regular part $\mathcal O^{reg}=E/Q$ by gluing back, at each
stratum, the orbifold cone $c(L_F)=N_F/F$ along the \emph{same} link $L_F$. The
orbifold fundamental group of a cone on a developable link with isotropy $F$ is
\[
   \pi_1^{orb}\big(c(L_F)\big)\ =\ F,
\]
and the inclusion $L_F\hookrightarrow c(L_F)$ induces on (orbifold) fundamental
groups exactly the surjection $\pi_1(L_F)\to F$ of \eqref{eq:linkgroup} (the cone
kills the base/fiber classes of the link and retains only the local isotropy $F$).
Hence the orbifold van Kampen theorem (Haefliger; for developable orbifolds it
agrees with the topological van Kampen of the orbifold's classifying space
$B\Gamma$ relative to $E\Gamma$) gives the parallel presentation
\begin{equation}\label{eq:vk-orb}
   \Gamma=\pi_1^{orb}(\mathcal O)\ \cong\ \pi_1(\mathcal O^{reg})\
   *_{\,\pi_1(L_F)\ (\text{all }F)}\ F .
\end{equation}
(That $\pi_1^{orb}(X/\Gamma)=\Gamma$ is standard: $X$ is contractible and
$\Gamma$ acts properly cocompactly, so $X$ is the orbifold universal cover and
$\Gamma$ its orbifold deck group.)

Now compare \eqref{eq:vk-M} and \eqref{eq:vk-orb}. They have the \emph{same}
free factor $\pi_1(\mathcal O^{reg})$ and the \emph{same} amalgamating link
groups $\pi_1(L_F)$ with the \emph{same} inclusions into $\pi_1(\mathcal
O^{reg})$. They differ only in the second amalgamated factor: $\pi_1(V_F)$ versus
$F$. But by the simply-connected normalization (M3, Lemma~\ref{lem:simpconn}) the
model satisfies $\pi_1(V_F)\cong F$, and \emph{this isomorphism is compatible
with the link inclusion}: the boundary monodromy
$\pi_1(L_F)\to\pi_1(V_F)\xrightarrow{\cong}F$ equals the structural surjection
$\pi_1(L_F)\to F$ of \eqref{eq:linkgroup} (by construction $W_F\to V_F$ is the
$F$-cover whose boundary restriction is $S(\nu)\to L_F$, so the monodromy on
$\pi_1(L_F)$ is precisely the covering monodromy of $S(\nu)\to L_F$, which is
\eqref{eq:linkgroup}). Therefore the inclusion-induced maps
$\pi_1(L_F)\to\pi_1(V_F)$ and $\pi_1(L_F)\to F$ are identified by the
isomorphism $\pi_1(V_F)\cong F$. An isomorphism of one amalgamated factor that
intertwines the amalgamating maps induces an isomorphism of the whole
amalgamated product (the amalgamated product is a colimit, functorial in the
factors over fixed amalgamating maps). Applying this at every stratum
simultaneously,
\[
   \pi_1(M)\ \cong\ \pi_1(\mathcal O^{reg})\,*_{\pi_1(L_F)}\,\pi_1(V_F)
   \ \cong\ \pi_1(\mathcal O^{reg})\,*_{\pi_1(L_F)}\,F\ \cong\ \Gamma .
\]
Thus $\pi_1(M)\cong\Gamma$, and the isomorphism is induced by the resolution map
$M\to\mathcal O$ (identity on $\mathcal O^{reg}$, collapse $V_F\to c(L_F)$ on the
models), so it is realized by the classifying map $c\colon M\to B\Gamma$ used
below.

\emph{The cover $\widetilde M$ is the simply connected universal cover.} Since
$M$ is a connected manifold (it is the quotient of the connected $\widehat M_0$,
which is connected because $M_0$ is and the resolution preserves connectedness as
above), it has a universal cover $\widetilde M\to M$, unique up to deck
isomorphism, with $\pi_1(\widetilde M)=1$ and deck group $\pi_1(M)\cong\Gamma$.
This universal cover is a smooth manifold (it pulls back the smooth atlas of $M$
along a local diffeomorphism) and carries the deck $\Gamma$-action freely,
properly discontinuously and cocompactly, with $\widetilde M/\Gamma=M$ compact
(the deck action of the fundamental group on the universal cover of a closed
manifold is always free, properly discontinuous and cocompact). In particular
$\widetilde M$ is \emph{simply connected}, hence genuinely the universal cover,
not merely some $\Gamma$-cover. (Equivalently: the cover of $M$ associated to the
normal subgroup $\ker(\pi_1(M)\xrightarrow{\cong}\Gamma)=1$ is the universal
cover, and its deck group is $\Gamma$.)

\emph{Rational acyclicity of $\widetilde M$.} We first show that the
classifying map $c\colon M\to B\Gamma$ (which induces the isomorphism
$\pi_1(M)\xrightarrow{\cong}\Gamma$ established above) is a rational homology
isomorphism, by descending the equivariant tube lemma to the Borel
constructions.

Since $Q$ acts \emph{freely} on $\widehat M_0$ with quotient $M$, the Borel
construction collapses:
\[
   H^Q_*(\widehat M_0;\mathbb{Q})
   = H_*(EQ\times_Q\widehat M_0;\mathbb{Q})
   = H_*(\widehat M_0/Q;\mathbb{Q}) = H_*(M;\mathbb{Q}),
\]
because $EQ\times_Q\widehat M_0\to\widehat M_0/Q=M$ is a fibration with
contractible fibers $EQ$. On the other side, $EQ\times_Q M_0=EQ\times_Q(X/\Gamma_0)$;
as $X$ is contractible and $\Gamma_0$ acts freely on it, $X/\Gamma_0\simeq B\Gamma_0$
$Q$-equivariantly, and the Borel construction of the extension
$1\to\Gamma_0\to\Gamma\to Q\to1$ gives a homotopy equivalence
$EQ\times_Q B\Gamma_0\simeq B\Gamma$ (equivalently
$EQ\times_Q M_0\simeq E\Gamma\times_\Gamma X\simeq B\Gamma$ since $X\simeq\ast$),
whence
\[
   H^Q_*(M_0;\mathbb{Q})=H_*(B\Gamma;\mathbb{Q})=H_*(\Gamma;\mathbb{Q}).
\]
By Proposition~\ref{prop:hatM} the collapse $r\colon\widehat M_0\to M_0$ induces
$H^Q_*(\widehat M_0;\mathbb{Q})\xrightarrow{\cong}H^Q_*(M_0;\mathbb{Q})$.
Combining the three displays,
\begin{equation}\label{eq:Mhomology}
   H_*(M;\mathbb{Q})\ \cong\ H_*(B\Gamma;\mathbb{Q})\ =\ H_*(\Gamma;\mathbb{Q}),
\end{equation}
and one checks directly that this isomorphism is induced by $c$: both $r$ and the
identifications above are compatible with the maps to $BQ$, and the composite
realizes $c_*$. Thus $c\colon M\to B\Gamma$ is a rational homology isomorphism
and a $\pi_1$-isomorphism. (Note this conclusion uses only that the $Q$-map $r$
is a \emph{rational} Borel-homology isomorphism and that $Q$ acts \emph{freely}
on $\widehat M_0$, which collapses $H^Q_*(\widehat M_0;\mathbb{Q})$ to
$H_*(M;\mathbb{Q})$; no integral or $\mathbb{Z}/2$-equivariant fundamental class
of $\widehat M_0$ or $M_0$ is transported, in accordance with
Remark~\ref{rem:equiv-transport}.)

Finally we deduce that the universal cover $\widetilde M$ is
$\mathbb{Q}$-acyclic, using the comparison theorem for Serre spectral sequences.
Consider the commutative ladder of fibrations over the common base $B\Gamma$:
\[
\begin{array}{ccccc}
\widetilde M & \longrightarrow & M & \xrightarrow{\ c\ } & B\Gamma\\[2pt]
\downarrow\widetilde c & & \downarrow c & & \big\|\\[2pt]
E\Gamma & \longrightarrow & B\Gamma & \xrightarrow{=} & B\Gamma,
\end{array}
\]
where the top row is the fibration with fiber the universal cover
$\widetilde M$ (its monodromy is the $\Gamma$-action on $H_*(\widetilde M)$), the
bottom row is the path fibration of $B\Gamma$ with contractible total space and
fiber $E\Gamma$, and $\widetilde c$ is the $\Gamma$-equivariant lift of $c$ to
universal covers (which exists and is unique up to homotopy because $c$ is a
$\pi_1$-isomorphism). Take homology with $\mathbb{Q}$ coefficients. The induced
map of homology Serre spectral sequences,
\[
E^2_{pq}=H_p\big(\Gamma;H_q(\widetilde M;\mathbb{Q})\big)\Rightarrow
H_{p+q}(M;\mathbb{Q}),\qquad
\bar E^2_{pq}=H_p\big(\Gamma;H_q(E\Gamma;\mathbb{Q})\big)\Rightarrow
H_{p+q}(B\Gamma;\mathbb{Q}),
\]
is an isomorphism on the base direction (the identity of $B\Gamma$) and, by
\eqref{eq:Mhomology}, an isomorphism on the abutments
$H_*(M;\mathbb{Q})\xrightarrow{\cong}H_*(B\Gamma;\mathbb{Q})$. By the
\emph{comparison theorem for spectral sequences of a map of fibrations with
identical base} (Zeeman comparison theorem; see e.g.\ McCleary,
\emph{A User's Guide to Spectral Sequences}, Thm.~3.26 and its homology dual),
if the map of abutments is an isomorphism and the map of base spectral sequences
is an isomorphism, then the map on the fiber is an isomorphism on homology:
\[
\widetilde c_*\colon H_*(\widetilde M;\mathbb{Q})\
\xrightarrow{\ \cong\ }\ H_*(E\Gamma;\mathbb{Q}).
\]
(The theorem applies: both are first-quadrant homology spectral sequences of
fibrations over the connected base $B\Gamma$, the comparison is induced by a map
of fibrations covering $\mathrm{id}_{B\Gamma}$, and on the fiber-degree-$0$ rows it
is the identity isomorphism $E^2_{p,0}=H_p(\Gamma;\mathbb{Q})=\bar E^2_{p,0}$
since both fibers are connected; equivalently, in the hyperhomology formulation
the map is the comparison
$\mathrm{Tor}^{\mathbb{Q}\Gamma}_*(\mathbb{Q},C_*(\widetilde M;\mathbb{Q}))\to
\mathrm{Tor}^{\mathbb{Q}\Gamma}_*(\mathbb{Q},C_*(E\Gamma;\mathbb{Q}))$ of two
bounded-below complexes of free $\mathbb{Q}\Gamma$-modules inducing an
isomorphism after $\mathbb{Q}\otimes_{\mathbb{Q}\Gamma}(-)$, whence the
comparison theorem yields an isomorphism on the input homologies $H_*(\widetilde
M;\mathbb{Q})$ and $H_*(E\Gamma;\mathbb{Q})$.) Since $E\Gamma$ is contractible,
$H_*(E\Gamma;\mathbb{Q})=H_*(\mathrm{pt};\mathbb{Q})$, so
\[
\widetilde H_*(\widetilde M;\mathbb{Q})=0,
\]
i.e.\ the universal cover $\widetilde M$ is $\mathbb{Q}$-acyclic.

\emph{Properness and cocompactness.} $\Gamma$ acts on $\widetilde M$ as the deck
group of the universal cover of the closed manifold $M$, hence freely, properly
discontinuously and cocompactly with $\widetilde M/\Gamma=M$ compact.

This establishes all assertions, and in particular the affirmative answer:
$\Gamma$ is realized as $\pi_1$ of the closed manifold $M$ with
$\mathbb{Q}$-acyclic universal cover.
\end{proof}

\begin{remark}[Scope: which lattices are covered unconditionally]\label{rem:scope}
The assembly above requires, at each isotropy stratum with stabilizer $F$, a
local model with boundary the resolved link \emph{on the nose}, which by
\S\ref{sec:induct} exists whenever the resolved link is null-bordant in oriented
bordism. This holds:
\begin{itemize}
\item \emph{unconditionally} when every stabilizer is $\mathbb{Z}/2$ (the minimal
$2$-torsion case): then each link is $\mathbb{RP}^{m-1}$ with $m$ even, which
bounds explicitly (Lemma~\ref{lem:RPodd}, Corollary~\ref{cor:Z2}). Hence for
every uniform lattice $\Gamma$ whose torsion is generated by involutions with
free slice representations (e.g.\ $\Gamma$ containing only $2$-torsion of order
exactly $2$ at the strata), Theorem~\ref{thm:main} is unconditional;
\item for general finite stabilizers, provided the resolved links bound
integrally (the stated bordism hypothesis of \S\ref{sec:induct}). In all cases
the bordism class is \emph{rationally} trivial, so after the $N$-fold device of
\S\ref{sec:induct} one obtains a global structure map that is a rational homology
equivalence; running the Borel-homology computation \eqref{eq:Mhomology} with this
rational equivalence still yields a closed manifold $M$ with $\pi_1=\Gamma$ and
$\mathbb{Q}$-acyclic universal cover, which is exactly the rational conclusion
asked for in Problem~7. The only nonformal input is thus the $2$-primary integral
bordism of the resolved links, isolated honestly in \S\ref{sec:induct}.
\end{itemize}
\end{remark}

\begin{remark}[Consistency with the obstructions]
$\widetilde M$ is $\mathbb{Q}$-acyclic but \emph{not} $\mathbb{Z}/2$-acyclic, and
indeed must not be: were it $\mathbb{Z}/2$-acyclic, Smith theory would force the
$2$-torsion subgroup of $\Gamma$ to fix a point of $\widetilde M$, contradicting
freeness of the $\Gamma$-action. The mod-$2$ homology of $\widetilde M$ is carried
exactly by the free models $W_F$ (Remark after Theorem~\ref{thm:atomic}), so
there is no conflict between the affirmative rational answer and the negative
integral/mod-$2$ statement.
\end{remark}

\section{Summary of what is established}\label{sec:summary}

\begin{enumerate}
\item The literal demand ``free cocompact action on a $\mathbb{Q}$-acyclic
manifold'' is impossible for $|F|>1$ (Lemma~\ref{lem:euler}); the correct local
object is a free $F$-manifold $W_F$ with $\mathbb{Q}$-acyclic \emph{quotient}
$V_F=W_F/F$ and boundary the link, which carries nontrivial mod-$2$ homology in
accordance with Smith theory.
\item The atomic existence theorem (Theorem~\ref{thm:atomic}): every rational
homology sphere $L^{m-1}$ ($m-1\ge5$) that bounds in oriented bordism bounds a
compact smooth $\mathbb{Q}$-acyclic manifold $V_F$; the free model $W_F$ is its
$F$-cover.
\item The minimal $2$-torsion case $F=\mathbb{Z}/2$ is settled
\emph{unconditionally} (Corollary~\ref{cor:Z2}): $\mathbb{RP}^{m-1}$
($m$ even, $\ge6$) bounds the explicit oriented manifold $D(\mathcal{O}(2))$
over $\mathbb{CP}^{(m-2)/2}$, which is surgered to $\mathbb{Q}$-acyclic.
\item General finite $2$-groups are handled by induction over isotropy strata
(Proposition~\ref{prop:induct}); the only nonformal input is integral
null-bordance of the resolved link, which holds rationally in all cases and
exactly for $\mathbb{Z}/2$. The rational statement is precisely what the global
resolution (issues \texttt{q7-5}, \texttt{q7-7}) consumes.
\item \textbf{Global resolution (issue \texttt{q7-5}, Theorem~\ref{thm:main}).}
Selberg's lemma presents $X/\Gamma=M_0/Q$ with $M_0=X/\Gamma_0$ closed
aspherical ($\Gamma_0$ torsion-free, free $X$-action) and $Q=\Gamma/\Gamma_0$
finite. The local models, grafted $Q$-equivariantly along the finitely many
isotropy strata (whose slice data coincide with the local $(F,V)$ by the Bochner
slice theorem, Lemma~\ref{lem:struct}), produce a closed smooth manifold
$\widehat M_0$ with a \emph{free} smooth $Q$-action and a $Q$-equivariant
rational homology equivalence $\widehat M_0\to M_0$ (equivariant tube lemma,
Lemma~\ref{lem:tube}, the precise form of \eqref{eq:RFNF}). Then
$M=\widehat M_0/Q$ is a closed (compact, \emph{boundaryless}) smooth manifold
with $\pi_1(M)\cong\Gamma$ (Seifert--van Kampen across the resolved strata,
comparing the pushout presentations \eqref{eq:vk-M} of $\pi_1(M)$ and
\eqref{eq:vk-orb} of $\pi_1^{orb}(\mathcal O)=\Gamma$; they agree because
$\pi_1(V_F)\cong F$ from (M3) intertwines the link monodromy
\eqref{eq:linkgroup}). Its \emph{simply connected} universal cover $\widetilde M$
is a smooth manifold on which $\Gamma$ acts freely, properly, cocompactly with
$\widetilde M/\Gamma=M$ and $\widetilde H_*(\widetilde M;\mathbb{Q})=0$.
Boundary-freeness is verified by collar-gluing (each frontier sphere bundle is
shared by exactly two pieces, $E$ and one resolved tube; deepest-first
processing keeps the frontiers disjoint closed manifolds), and $\widetilde M$ is
genuinely the universal cover because $\ker(\pi_1(M)\xrightarrow{\cong}\Gamma)=1$
(Proposition~\ref{prop:hatM}, Theorem~\ref{thm:main}). This gives the affirmative answer to
Problem~7 for every uniform lattice $\Gamma$, including those with $2$-torsion;
consistently, $\widetilde M$ is not $\mathbb{Z}/2$-acyclic (Smith theory).
\end{enumerate}
\hfill$\qed$

\end{document}
